\chapter{Economic implications: a deliberately cautious estimate}
\label{ch:economics}

\begin{quote}\itshape
This chapter in brief: what the measured results could be worth to the three parties who pay
for engine health monitoring --- maintenance shops, airlines and original-equipment
manufacturers (OEMs) --- for one concrete
operator, with every cost input ranged and cited, every dollar traced to a measured
mechanism, the uncertainty propagated rather than hidden, and an explicit section on why the
number might be zero. It leads with the caveats on purpose.
\end{quote}

\keyidea{Only one mechanism is monetized --- less-biased prognosis converting unscheduled
engine removals into scheduled ones --- because it is the only one the results directly
support; the estimate is anchored to industry rates (not the synthetic fleet's optimistic
raw numbers), propagated through Monte Carlo, and reported as a wide interval with its
downside.}

\section{The rule: no dollar without a measured mechanism}
\label{sec:econ-rule}

It is easy to manufacture a large return on investment (ROI) for any monitoring technology; it is honest to
manufacture none. This chapter monetizes exactly one mechanism, the one \cref{ch:results}
measured: the tuned AI prognosis is markedly less biased and tighter than the tuned
traditional extrapolation at end of life (RUL RMSE $858$ vs $1\,981$ cycles at $90\,\%$
life). A less-optimistic remaining-life estimate lets an operator \emph{slot} a shop visit
before the engine fails on the wing --- converting a fraction of expensive
\emph{unscheduled} removals (AOG, spare-engine lease, expedited logistics, higher risk of
secondary damage) into cheaper \emph{scheduled} ones. Everything else an enthusiast might
count --- fuel saved by earlier fault clearing, smoother shop capacity, the OEM value of
better sensors --- is discussed but \emph{not} put in the total, because our benchmark does
not measure it directly. All inputs live in \code{conf/econ\_assumptions.yaml} as ranges
with sources or explicit \textsc{assumption} flags; the arithmetic is
\code{scripts/econ\_impact.py} (norm N4).

\section{The concrete case}
\label{sec:econ-case}

A mid-size 737NG operator: 50 aircraft, hence 100 CFM56-7B engines, at a public utilization
of 2\,000--3\,000 cycles per engine per year. A mature CFM56-7B shop visit is taken as
\$3--5\,M (public maintenance-repair-overhaul (MRO) commentary, wide because workscope varies); an unscheduled removal
costs $1.3$--$2.5\times$ that; the monitoring programme itself costs \$0.1--1\,M per year to
run (software, integration, analyst time, false-alarm handling), charged against any saving.
Every one of these is an uncertain range, and the Monte Carlo (20\,000 trials) draws all of
them jointly.

\paragraph{Honesty about the grounding.} On the synthetic fleet, a naive rule that removes
purely when the point RUL crosses the lead time would make the traditional method
$40\,\%$ unscheduled and the AI $0\,\%$ --- figures that \emph{overstate} reality, because
real trend monitoring uses conservative margins and more signals than our benchmark's
cockpit RUL. So those numbers are used only to establish the \emph{direction and rough size}
of the improvement; the headline \textbf{anchors the absolute unscheduled rate to the
industry band ($10$--$30\,\%$)} and lets better prognosis convert a deliberately
conservative, wide $15$--$60\,\%$ of those removals. This anchoring is the single most
important honesty decision in the chapter.

\section{The mechanism, measured: unscheduled$\,\to\,$scheduled conversion}
\label{sec:ops-conversion}

The naive $40\,\%$-vs-$0\,\%$ figure above is too crude to build on. Before committing any
dollar to the single monetized mechanism, we measure it properly (tag \texttt{prereg-v12}):
\emph{what fraction of run-to-failure removals does each approach actually convert into planned
ones}, as a function of how much warning the logistics chain needs. This turns the abstract RUL
error of \cref{ch:results} into the operational number the estimate below consumes.

\keyidea{An engine reaches its booked shop slot without failing first only if the method's RUL
\emph{over-prediction} stays within the logistics horizon $L$; over-predict by more than $L$ and
the removal is unscheduled. Netting out engines pulled wastefully early (the prognosis analog of
a false alarm), the AI converts more at every life stage --- and, tellingly, sits near the
irreducible F11 ceiling early in life but far below it late, exactly where a better prognostic
still has room.}

\paragraph{The operational rule.} At an inspection taken at life fraction $f$, an engine's true
remaining life is $R$ and the method predicts $\hat R = R + e$, with $e$ the signed RUL error
already recorded in \cref{ch:results} ($e>0$ is over-prediction). The operator books a removal a
horizon $L$ cycles before the predicted end. The engine survives to that slot --- a
\textbf{converted, scheduled} removal --- iff $e \le L$; if $e > L$ it fails first
(\textbf{unscheduled}); if it grossly under-predicts ($e < -W$) the engine is pulled far too
early (a \textbf{wasteful} removal, the false-alarm analog). The honest headline is the
\textbf{net} rate, $-W \le e \le L$, which credits only removals that are both on time and not
wasteful. Nominal $L=400$, $W=800$ cycles; the horizon is swept, not chosen.

\begin{figure}[htbp]
  \centering
  \includegraphics[width=\textwidth]{ops_conversion}
  \caption{Net unscheduled$\to$scheduled conversion versus logistics horizon $L$, per method,
  at three inspection points. Baseline (no monitoring) is $0\,\%$. The grey ceiling is what an
  \emph{unbiased} predictor limited only by the F11 aleatoric floor
  (\cref{sec:bt-prognostic-floor}) could achieve --- the shaded band above it is physically
  unreachable. Early in life the AI already hugs that ceiling; late in life a wide reachable gap
  remains.}
  \label{fig:ops-conversion}
\end{figure}

\paragraph{What the measurement shows.} Three results, frozen as hypotheses before the pass:

\begin{itemize}
  \item \textbf{H-OPS.1 confirmed --- the AI converts more, net.} At the nominal horizon the AI
    converts $35/50/45\,\%$ of removals at $50/70/90\,\%$ life, against the traditional
    $15/0/25\,\%$ --- a strict win at every stage. The traditional stack's optimistic,
    high-variance RUL (\cref{ch:trad}) over-predicts past the horizon too often to plan around.
  \item \textbf{H-OPS.2 confirmed --- the ceiling binds early, headroom opens late.} The gap
    between the physical ceiling and the AI's achieved conversion is essentially zero at
    mid-life (the AI hugs the $\sim\!65\,\%$ floor-limited ceiling) but $52$ points at
    $90\,\%$ life (ceiling $97\,\%$, AI $45\,\%$). This is the operational face of the
    prognostic floor: early conversion is capped by irreducible randomness, so buying a better
    prognostic pays off \emph{late}, not early --- the same verdict \cref{sec:bt-prognostic-floor}
    reached in RMSE terms, now in removals.
  \item \textbf{H-OPS.3 confirmed --- the honest accounting is sobering but survives.} Netting
    out wasteful early removals nearly halves the AI's raw conversion at mid-life ($65\to35\,\%$):
    part of the apparent gain is just aggressive early pulls. Even so, the net AI rate still
    beats the net traditional rate. The advantage is real; it is smaller than the gross number
    would suggest.
\end{itemize}

These synthetic net rates ($35$--$50\,\%$) sit \emph{inside} the deliberately conservative
$15$--$60\,\%$ band the estimate anchors to, which is why the chapter uses the anchored band
rather than these point values: with $n=20$ test engines the bootstrap intervals are wide, and
the synthetic fleet cannot stand in for a real operator's margins and extra signals. The
measurement earns the \emph{direction and rough size} of the mechanism; the estimate keeps the
absolute rate anchored to industry.

\section{The estimate, with its uncertainty}
\label{sec:econ-result}

With the mechanism fixed and the inputs ranged, the Monte Carlo turns them into a
distribution of annual savings rather than a single number --- and the distribution, not its
midpoint, is the honest answer.

\begin{table}[htbp]
  \centering\small
  \caption{Annual economic impact for the 100-engine operator, from the mechanism our
  results support. Monte-Carlo percentiles over all ranged inputs. Auto-generated.}
  \label{tab:econ}
  % AUTOGENERATED by scripts/make_report_assets.py — do NOT edit by hand
\begin{tabular}{lr}
\toprule
Annual fleet impact (100 engines) & Value \\
\midrule
  p10 (conservative) & \$0.9\,M \\
  p50 (median) & \$3.2\,M \\
  p90 (optimistic) & \$8.0\,M \\
  per engine (p50) & \$32\,k/yr \\
  probability net-negative & 1.0\,\% \\
\bottomrule
\end{tabular}

\end{table}

\Cref{tab:econ,fig:econ}: the median annual saving is about \textbf{\$3\,M for the
100-engine fleet --- roughly \$30\,k per engine per year} --- but the honest content is the
spread, not the median. The conservative decile (p10) is still positive at \$0.9\,M, the
optimistic decile (p90) reaches \$8\,M, and the estimate is net-negative in about $1\,\%$ of
parameter draws. That $1\,\%$ understates the true risk, because it counts only
\emph{parametric} uncertainty; the structural risks of \cref{sec:econ-worthless} are not in
it. The sensitivity panel shows the outcome is driven most by the unscheduled-cost premium
and the conversion fraction --- precisely the two softest assumptions --- so a reader who
distrusts those should read toward p10, not p50.

\begin{figure}[htbp]
  \centering
  \includegraphics[width=\textwidth]{econ_impact}
  \caption{Left: Monte-Carlo distribution of annual fleet savings (red line at zero, dashed
  at the median). Right: sensitivity of the outcome to each input --- the two softest
  assumptions dominate. Auto-generated.}
  \label{fig:econ}
\end{figure}

\section{The three parties}
\label{sec:econ-strata}

The saving above is an operator's; but engine health monitoring is paid for by three parties
with different stakes, and the results speak to each differently.

\begin{description}
  \item[Airlines (operators).] The mechanism above is theirs: fewer AOG events and plannable
    shop slots. The quantified \$0.9--8\,M/year is an operator number. A second, unmonetized
    operator benefit is fuel: clearing a fault (or scheduling a wash) earlier trims the
    fuel-flow penalty of degradation, but our benchmark does not measure fleet fuel, so it is
    left out of the total rather than guessed.
  \item[Maintenance shops.] The value is capacity smoothing: unscheduled removals
    arrive at random and force expensive expediting, while scheduled ones fill planned slots.
    Part of this is already inside the unscheduled-cost premium; the residual --- better
    workforce and slot utilization --- is real but operator-specific and left qualitative.
    The determinability certificate (\cref{ch:breakthrough}) additionally enables
    \emph{condition-based workscoping}: knowing which module is genuinely degraded (when it
    is identifiable) avoids opening what does not need opening.
  \item[OEMs and sensor suppliers.] Here the project's diagnostic results become a design
    input rather than a saving. The certificate quantifies exactly which added sensor rescues
    which fault --- station probes make HPC-efficiency faults $45\times$ more identifiable
    (\cref{sec:bt-honesty}), fan efficiency almost not at all --- turning ``instrument the
    engine better'' into a ranked, physics-justified specification. We deliberately attach no
    dollar figure to this: its value depends on an OEM's product strategy, which our data
    cannot see.
\end{description}

\section{Why this could be worth nothing}
\label{sec:econ-worthless}

The estimate rests on the mechanism actually materializing in service, and several structural
risks --- none in the Monte Carlo --- could drive it to zero or below:
\begin{itemize}
  \item \textbf{The isolation wall (H2) is unbroken.} The AI does \emph{not} improve fault
    isolation (\cref{sec:res-h2}); its value here is purely prognostic. An operator hoping
    for better \emph{diagnosis} from the same investment would be disappointed.
  \item \textbf{False alarms erode trust and money.} A prognosis that cries wolf triggers
    premature removals --- the opposite of the saving --- and maintenance controllers stop
    believing it. The integration cost line captures some of this, but a badly-tuned
    deployment could be net-negative well beyond the modelled $1\,\%$.
  \item \textbf{Sim-to-real gap.} Our RUL advantage is measured on a synthetic fleet; the
    C-MAPSS check (\cref{sec:sim-to-real}) preserved the \emph{ranking} but not necessarily
    the magnitude. Real degradation is messier, and the conversion fraction could sit at the
    bottom of its range or below.
  \item \textbf{Incumbent practice is already good.} Mature operators catch most failures
    today with conservative margins; the marginal removals an AI actually converts may be
    fewer than even our anchored band assumes.
\end{itemize}

\section{What the results support, and what they do not}
\label{sec:econ-support}

Supported by measurement: a less-biased end-of-life RUL exists (H3, confirmed and externally
ranked), and its \emph{direction} of economic effect --- fewer unscheduled removals --- is
sound. \emph{Not} established by this work: the absolute magnitude in a real fleet, which
depends on deployment quality, incumbent practice and real degradation, none observable here.
The honest summary is therefore a conditional one: \textbf{if} the prognosis advantage
transfers even partially to service and \textbf{if} the programme is competently integrated,
a 100-engine 737NG operator can plausibly save low single-digit millions per year, with a
real chance of less and a structural chance of nothing --- and the same analysis says the
larger, unquantified prize is not prognosis at all but the diagnostic instrumentation the
certificate would justify buying.
